% Appendix: its own sections, so the mini-frame headline shows only these.
\section[Templates]{Templates}

% ------------------------------------------------------------- AGENTS.md
\begin{frame}[fragile]{A1\enspace Template: \texttt{AGENTS.md}}
  \begin{columns}[T,onlytextwidth]
    \begin{column}{0.62\textwidth}
\begin{codefile}[basicstyle=\ttfamily\fontsize{7}{8.4}\selectfont]{AGENTS.md}
# <project>: notes for coding agents
## Setup and checks
- Environment: `<create and activate command>`
- Fast check (CPU, synthetic data, under a minute): `<command>`
- Lint and types: `<command>`
## Where things are
- Vocabulary: `GLOSSARY.md`; decisions: `docs/adr/`
## Conventions the code does not show
- <units, coordinate frames, file formats, splitting rules>
## Boundaries
- Read-only: `data/raw/`, <other protected paths>
- Ask before: new dependencies, deleting files, cluster jobs
- Never: commit secrets or data; push to `main`
- Never edit a test just to make it pass
## When you finish
- Changed files, commands run, assumptions, open questions
\end{codefile}
    \end{column}
    \begin{column}{0.35\textwidth}
      \footnotesize
      \begin{itemize}\setlength\itemsep{1.4mm}
        \item each line prevents a mistake you have seen
        \item point to documents; do not copy them
        \item update it in the commit that makes it stale
        \item Codex: nested files per sub-package; an override replaces one folder's file
        \item other tools read other names, e.g.\ \fn{CLAUDE.md}
      \end{itemize}
    \end{column}
  \end{columns}
\end{frame}

% --------------------------------------- glossary, coding standards, ADR
\begin{frame}[fragile,source={Nygard (2011) · Pocock (2026), domain-modeling ADR format; code-review}]{A2\enspace Templates: glossary, coding standards and ADR}
  \begin{columns}[T,onlytextwidth]
    \begin{column}{0.48\textwidth}
\begin{codefile}[basicstyle=\ttfamily\fontsize{7}{8.4}\selectfont]{GLOSSARY.md}
# Glossary
<Term>: one-sentence definition (`code_name`).
  Relates to: <other terms and how>.
  Avoid: <synonyms we do not use>.
<Term>: ...
\end{codefile}
      \vskip2mm
\begin{codefile}[basicstyle=\ttfamily\fontsize{7}{8.4}\selectfont]{CODING\_STANDARDS.md}
# Coding standards
- Randomness: <how seeds are passed>
- Errors: <what public functions raise>
- Tests: <fixtures, seams, expected values>
- <a judgement call a reviewer can check>
\end{codefile}
    \end{column}
    \begin{column}{0.49\textwidth}
\begin{codefile}[basicstyle=\ttfamily\fontsize{7}{8.4}\selectfont]{docs/adr/NNNN-short-title.md}
# ADR-NNNN: <short noun phrase>
Status: proposed | accepted |
        superseded by ADR-NNNN
## Context
Forces at play: data, constraints, question.
## Decision
What we will do, in full sentences.
## Considered options
Alternatives and why they lost.
## Consequences
What follows, including the costs.
\end{codefile}
    \end{column}
  \end{columns}
\end{frame}

% ------------------------------------------------------- task specification
\begin{frame}[fragile]{A3\enspace Template: bounded task specification}
  \begin{columns}[T,onlytextwidth]
    \begin{column}{0.62\textwidth}
\begin{codefile}[basicstyle=\ttfamily\fontsize{7}{8.1}\selectfont]{spec.md}
# Task: <one-line objective>
## Objective
What is true afterwards? Name modules, functions, options.
## Context
Glossary terms, relevant ADRs, input and output formats.
## Constraints
Files not to touch; dependencies, data and compute limits.
## Acceptance criteria
- Observable behaviour, including failure cases
- Determinism and reproducibility requirements
## Test seams
Public interfaces the tests use; nothing else is tested.
## Verification
- Command to run and the expected result
- Report: changed files, test output, assumptions
## Out of scope
What must not change in this task.
\end{codefile}
    \end{column}
    \begin{column}{0.35\textwidth}
      \footnotesize
      \begin{itemize}\setlength\itemsep{1.4mm}
        \item acceptance criteria as observable behaviour
        \item at least one failure case, e.g.\ fewer specimens than folds
        \item the test seams; the first test approved before the code
        \item the verification command: the agent runs it
        \item what is out of scope: agents tidy up beyond the task
        \item in a prepared repository, context and constraints are pointers to files
      \end{itemize}
    \end{column}
  \end{columns}
\end{frame}

\section[Tool details]{Tool details}

% --------------------------------------------------------------------- uv
% A4 and A5 are kept for questions after slide 11; their notes say so instead of giving a time
\begin{frame}[fragile,source={Astral (2026), uv documentation: projects, running commands, lockfile; uv README and its speed benchmark}]{A4\enspace The upgrade: uv}
  \begin{columns}[T,onlytextwidth]
    \begin{column}{0.49\textwidth}
      {\footnotesize\bfseries A uv project\par}\vskip1.2mm
      \begin{foldertree}[note at=29.5]
        \tdir{0}{sem-seg/}{all committed except \fn{.venv/}}
        \tfile{1}{pyproject.toml}{what the project needs}
        \tfile{1}{uv.lock}{exact versions, every platform}
        \tfile{1}{.python-version}{the Python version}
        \tdir{1}{.venv/}{rebuilt from \fn{uv.lock}}
      \end{foldertree}
      \footnotesize
      \begin{itemize}\setlength\itemsep{1.2mm}
        \item one tool instead of pip, venv, pip-tools, pipx and pyenv
        \item much faster than pip
        \item \fn{uv run}: nothing to activate; the environment is synced first
      \end{itemize}
    \end{column}
    \begin{column}{0.48\textwidth}
\begin{terminal}[basicstyle=\ttfamily\fontsize{7}{8.6}\selectfont]{terminal}
$ uv init sem-seg            # new project
$ uv python pin 3.12         # .python-version
$ uv add numpy pandas torch  # updates uv.lock
$ uv add --dev pytest ruff mypy
$ uv run pytest -q           # sync, then run
$ uv sync                    # rebuild .venv
\end{terminal}
    \end{column}
  \end{columns}
  \note{\textbf{If asked}\par
    uv, from Astral, the makers of Ruff, is one tool for what pip, venv, pip-tools, pipx and
      pyenv do separately.\par
    A project is a pyproject.toml with what the project needs and a uv.lock with the exact
      version of every package for every platform; commit both.\par
    uv builds the environment in .venv inside the project; it is not committed and can be
      rebuilt at any time.\par
    uv add adds a package and updates the lockfile; tools for development, such as pytest, Ruff
      and mypy, go into their own group.\par
    uv run checks that lockfile and environment are in sync and then runs the command, so there
      is nothing to activate, which suits an agent whose commands each start a fresh shell.\par
    Astral reports 10 to 100 times the speed of pip.\par
    The current version is 0.13; uv has not reached 1.0, but Astral treats it as stable.}
\end{frame}

% ------------------------------------------------------------ conda or uv
\begin{frame}[source={conda (2026) · Astral (2026), uv: PyTorch integration · prefix.dev (2026), pixi}]{A5\enspace conda or uv?}
  % one question decides; the rest of the comparison is what slide 11 and A4 show
  \begin{tikzpicture}[x=1mm,y=1mm,
      ask/.style={human,text width=60mm,inner sep=1.6mm,font=\footnotesize,align=center},
      tool/.style={agent,anchor=north,text width=58mm,inner sep=1.8mm,font=\footnotesize,align=left},
      yn/.style={lbl,font=\scriptsize\bfseries}]
    \def\fact{\ttfamily\fontsize{7}{8.4}\selectfont}
    \node[ask] (q) at (72,-4.6) {\textbf{needs non-Python packages?}\\[0.3mm]
      \scriptsize CUDA toolkits, GDAL, compilers};
    \node[tool] (c) at (34,-17) {{\color{udsblue}\bfseries conda}\enspace from conda-forge, or pixi\\[1mm]
      \fact environment.yml + lockfile\\ conda run -n sem-seg \dots};
    \node[tool] (u) at (110,-17) {{\color{udsblue}\bfseries uv}\enspace Python packages from PyPI\\[1mm]
      \fact pyproject.toml + uv.lock\\ uv run \dots};
    \draw[flow] (q.west) -| node[yn,pos=0.75,left] {yes} (c.north);
    \draw[flow] (q.east) -| node[yn,pos=0.75,right] {no} (u.north);
    % the one extra step for GPU work with uv
    \node[ask,anchor=north,text width=50mm] (g) at (110,-35.4) {\textbf{GPU PyTorch on Windows?}};
    \draw[flow] (u.south) -- (g.north);
    \node[agent,anchor=north,text width=50mm,inner sep=1.6mm,font=\footnotesize,align=center] (i) at (110,-45.2)
      {PyTorch's index in \fn{pyproject.toml}\\ \scriptsize PyPI's Windows wheels are CPU-only};
    \draw[flow] (g.south) -- node[yn,right] {yes} (i.north);
    % common ground
    \node[anchor=north west,inner sep=0pt,font=\scriptsize,text=udsdark,text width=64mm,align=left] at (2,-35.6)
      {\textbf{either way}\\ one project file defines the environment\\ exact versions in a lockfile\\
       one command runs the code\\ about five commands to learn};
  \end{tikzpicture}
  \note{\textbf{If asked}\par
    One question decides: does the project need packages that are not Python?\par
    conda installs packages of any language from conda-forge, including CUDA toolkits, GDAL or
      compilers; uv installs Python packages from PyPI and nothing else.\par
    Both can now pin exact versions: uv with its lockfile, conda with its own lockfile since
      version 26.5 or with conda-lock.\par
    conda keeps environments in a central folder, uv in the project.\par
    For GPU work, conda-forge has CUDA builds, while PyPI's PyTorch wheels for Windows are
      CPU-only, so a uv project needs PyTorch's own index in its pyproject.toml, as the uv guide
      shows.\par
    The learning curve is small: a handful of commands, and the habit that the project file, not
      the shell, defines the environment.\par
    If a project needs non-Python packages, stay with conda, or look at pixi, which installs
      conda-forge packages with a lockfile and uses uv for PyPI packages.}
\end{frame}

% ------------------------------------------------------- Codex discovery
\begin{frame}[source={OpenAI (2026), Custom instructions with AGENTS.md · Codex configuration reference}]
  {A6\enspace Codex: how instruction files are found}
  \begin{tikzpicture}[x=1mm,y=1mm,
      file/.style={artifact,minimum width=33mm,minimum height=5.2mm,inner sep=0.6mm,font=\ttfamily\scriptsize},
      dir/.style={box,minimum width=40mm,minimum height=5.2mm,inner sep=0.6mm,font=\ttfamily\scriptsize}]
    % ---- left: one folder
    \node[anchor=north west,inner sep=0pt,font=\small\bfseries] at (0,0) {Within one folder: replacement};
    \node[file,anchor=north west] (o1) at (2,-6.5) {AGENTS.override.md};
    \node[file,anchor=north west] (o2) at (2,-14)  {AGENTS.md};
    \node[file,anchor=north west] (o3) at (2,-21.5){fallback names};
    \foreach \a/\b in {o1/o2,o2/o3} \draw[flow] (\a.south) -- (\b.north);
    \draw[udsgreymid,line width=0.5pt] ([xshift=2mm]o1.north east) -- ++(1.2,0) |- ([xshift=2mm]o3.south east);
    \node[anchor=west,text width=28mm,align=left,font=\scriptsize,inner sep=0pt,
      execute at begin node={\hyphenpenalty=10000}] at ([xshift=4.5mm]o2.east)
      {checked top to bottom; the first non-empty file is used};
    \node[anchor=north west,text width=64mm,font=\footnotesize,inner sep=0pt] at (0,-31)
      {At most \textbf{one file per folder}: an override \emph{replaces} the \fn{AGENTS.md} beside it.};
    \draw[udsgreylight,line width=0.4pt] (69,0) -- (69,-41);
    % ---- right: chain of folders
    \node[anchor=north west,inner sep=0pt,font=\small\bfseries] at (73,0) {Across folders: layering};
    \foreach \t [count=\i from 0] in
        {\textasciitilde/.codex/\ (global),repository root (.git),src/,src/sem\_seg/,src/sem\_seg/data/\ \ <- cwd} {
      \node[dir,anchor=north west] (d\i) at (74,-6-\i*6.7) {\t};
    }
    \draw[udsred,line width=0.9pt] (d4.north west) rectangle (d4.south east);
    \foreach \i [evaluate=\i as \j using int(\i+1)] in {0,...,3} \draw[flow] (d\i.south) -- (d\j.north);
    \draw[udsblue,line width=0.6pt,-{Stealth[length=1.7mm]}] (119,-6.5) -- (119,-37.3);
    \node[anchor=west,text width=22.5mm,align=left,font=\scriptsize,inner sep=0pt] at (120.5,-22)
      {concatenated top to bottom; later files take \textbf{precedence}};
    \node[anchor=north west,font=\footnotesize,inner sep=0pt] at (73,-40.5)
      {Folders below or beside \texttt{cwd}: \textbf{not} loaded.};
    % ---- caveats
    \foreach \t [count=\i from 0] in
      {{Built \textbf{once per run}: restart Codex after editing an instruction file.},
       {\fn{project\_doc\_max\_bytes} (default 32\,KiB) limits instruction \emph{text}; it is not the model's context window.},
       {Check what loaded: ask Codex to summarize its current instructions (command in A7).}} {
      \node[box,fill=darktint,draw=none,anchor=north west,text width=43mm,align=left,
        font=\scriptsize,minimum height=9.5mm] at (\i*49,-44.5) {\t};
    }
  \end{tikzpicture}
\end{frame}

% ---------------------------------------------------------- Codex details
\begin{frame}[source={OpenAI (2026), Codex documentation and source code: AGENTS.md, configuration, sandboxing, approvals, skills}]{A7\enspace Codex: files, limits and controls}
  \scriptsize\renewcommand\arraystretch{1.2}\setlength\tabcolsep{3pt}
  \begin{tabularx}{\textwidth}{@{}>{\bfseries}L{19mm}X@{}}
    \toprule
    Global file  & In \fn{\$CODEX\_HOME} (default \fn{\textasciitilde/.codex}): \fn{AGENTS.override.md} if present and
                   non-empty, otherwise \fn{AGENTS.md}; only the first non-empty file is used.\\
    Project root & Found by walking up from the working directory to a root marker (default \fn{.git};
                   \fn{project\_root\_markers}). Without a root, only the working directory is checked.\\
    Per folder   & \fn{AGENTS.override.md}, then \fn{AGENTS.md}, then names in \fn{project\_doc\_fallback\_filenames};
                   at most one file per folder.\\
    Order        & Concatenated from the root down to the working directory, joined by blank lines; files
                   closer to the working directory come later and override earlier guidance.\\
    Size         & \fn{project\_doc\_max\_bytes}, 32\,KiB by default; Codex stops adding files at the limit.\\
    Timing       & Built once per run, usually once per interactive session; restart after editing.\\
    Checks       & \fn{codex --ask-for-approval never "Summarize the current instructions."};
                   \fn{codex --cd <folder> ...} for nested folders; logs with \fn{-c log\_dir=./.codex-log}.\\
    Sandbox      & \fn{read-only}, \fn{workspace-write} (the usual local mode), \fn{danger-full-access}. In
                   \fn{workspace-write} the network is off unless enabled; \fn{.git}, \fn{.agents}, \fn{.codex},
                   \fn{.aws} stay read-only.\\
    Approvals    & \fn{on-request}, \fn{never} or granular rules; \fn{untrusted} is retired; \fn{on-failure} is
                   read as \fn{on-request}.\\
    Skills       & Folders with \fn{SKILL.md} in \fn{.agents/skills} (repository) or \fn{\$HOME/.agents/skills};
                   invoked with \fn{\$name} or \fn{/skills}, or implicitly when a task matches.\\
    \bottomrule
  \end{tabularx}
\end{frame}

% ------------------------------------------------- Pocock: video vs current
\begin{frame}[source={Pocock (2026), video of 16 Jul 2026 · mattpocock/skills v1.3.1 · aihero.dev/skills}]{A8\enspace Pocock's workflow, then and now}
  \scriptsize\renewcommand\arraystretch{1.28}\setlength\tabcolsep{3pt}
  \begin{tabularx}{\textwidth}{@{}>{\bfseries}L{22mm}L{55mm}X@{}}
    \toprule
    & \bfseries Video (16 July 2026, between v1.1 and v1.2) & \bfseries Documentation, v1.3.1 (October 2026)\\
    \midrule
    Domain file  & \fn{CONTEXT.md} plus ADRs & \fn{GLOSSARY.md} plus ADRs; the skills read only \fn{GLOSSARY.md}\\
    Main flow    & \fn{/grill-with-docs}, then \fn{/implement} directly, or \fn{/to-spec}, \fn{/to-tickets},
                   \fn{/implement} (which runs \fn{/code-review}) &
                   \fn{/grill-with-docs}, \fn{/to-spec}, \fn{/to-tickets}, \fn{/implement} (which drives \fn{/tdd}
                   and runs \fn{/code-review}); \fn{/implement-spec} for a whole spec\\
    Grilling     & one question at a time & numbered rounds, each question with a recommended answer\\
    Local tickets& one \fn{tickets.md} & one file per ticket under \fn{.scratch/<feature>/issues/}\\
    Context      & ``smart zone'' of about 140k tokens (heuristic); clear before each ticket &
                   about 150k tokens (heuristic); one session for clarifying to ticketing, a fresh one per ticket\\
    Testing      & not shown & \fn{/tdd}: red to green at agreed seams; refactoring moved to \fn{/code-review}\\
    Installation & \fn{npx skills@latest add mattpocock/skills} & Claude Code plugin recommended; \fn{npx} for other agents\\
    \bottomrule
  \end{tabularx}
\end{frame}

\section[Sources]{Sources}

% -------------------------------------------------------------- references
\begin{frame}{A9\enspace References: tools and workflow}
  \scriptsize\setlength\parskip{0.9mm}\vskip-\parskip
  OpenAI (2026). \emph{Custom instructions with AGENTS.md}.
  \url{https://learn.chatgpt.com/docs/agent-configuration/agents-md}, accessed 19~Oct~2026.\par
  OpenAI (2026). Codex documentation: \emph{Configuration reference}, \emph{Sandboxing}, \emph{Agent approvals and
  security}, \emph{Non-interactive mode}, \emph{Build skills}. \url{https://learn.chatgpt.com/docs}, accessed 19--22~Oct~2026;
  source code, \url{https://github.com/openai/codex}.\par
  OpenAI (2026). \emph{Running Codex safely at OpenAI}. 8~May~2026. \url{https://openai.com/index/running-codex-safely/}\par
  AGENTS.md: a simple, open format for guiding coding agents. \url{https://agents.md/}, accessed 19~Oct~2026.\par
  Agent Skills format. \url{https://agentskills.io/}, accessed 19~Oct~2026.\par
  Anthropic (2026). Claude Code documentation: \emph{Best practices}, \emph{Memory}, \emph{Checkpointing}, \emph{Tools
  reference}, \emph{Worktrees}, \emph{Desktop}, \emph{CLI reference}. \url{https://code.claude.com/docs},
  accessed 20~Oct~2026.\par
  Microsoft (2026). VS Code documentation: \emph{Source control}, \emph{Python environments}, \emph{Editing Python}.
  \url{https://code.visualstudio.com/docs}, accessed 20~Oct~2026.\par
  M.~Pocock (2026). \emph{mattpocock/skills}, v1.3.1, commit 49dd158 (9~Oct~2026), MIT licence.
  \url{https://github.com/mattpocock/skills}\par
  M.~Pocock (2026). \emph{AI skills for real engineers}, \url{https://www.aihero.dev/skills}, and \emph{v1.3: /implement-spec,
  /pr, /retro, and GLOSSARY.md}, 5~Oct~2026,
  \url{https://www.aihero.dev/skills-changelog-v13-implement-spec-pr-retro-and-glossary-md}\par
  M.~Pocock (2026). \emph{mattpocock/skills: A complete AI coding workflow, end-to-end}. YouTube, 16~Jul~2026.
  \url{https://www.youtube.com/watch?v=M6mYodf0dJM}\par
\end{frame}

\begin{frame}{A10\enspace References: git, environments and checks}
  \scriptsize\setlength\parskip{0.9mm}\vskip-\parskip
  S.~Chacon, B.~Straub (2014). \emph{Pro Git}, 2nd ed. Apress. \url{https://git-scm.com/book}; git documentation,
  \emph{git-worktree}, \url{https://git-scm.com/docs/git-worktree}, accessed 20--21~Oct~2026.\par
  Layout of the worktree diagram, accessed 21~Oct~2026: Anthropic Academy (2026), figure of worktree branches
  merged into main, \url{https://academy.claude.com/assets/media/536e6bef11940caea62b1d41d10070eb1e2a220811cb0ff388da1af20b4d2c79.png};
  DevToolLab (2026), worktrees guide, \url{https://devtoollab.com/blog/claude-code-git-worktrees-parallel-agents-guide}.\par
  conda (2026). Documentation, v26.9: \emph{Managing environments}, \emph{conda run}, \emph{New features}.
  \url{https://docs.conda.io}, accessed 20~Oct~2026.\par
  Anaconda (2025). \emph{Terms of service}, 15~Jul~2025; \emph{Academic policy}, 14~Jul~2025.
  \url{https://www.anaconda.com/legal/terms/terms-of-service}\par
  conda-forge (2026). \emph{Miniforge}. \url{https://github.com/conda-forge/miniforge}; prefix.dev (2026). \emph{pixi
  documentation}, v0.82. \url{https://pixi.prefix.dev}\par
  Astral (2026). \emph{uv documentation}, v0.13: projects, running commands, PyTorch integration.
  \url{https://docs.astral.sh/uv/}, accessed 20~Oct~2026; README with its speed benchmark, \url{https://github.com/astral-sh/uv}.\par
  Astral (2026). \emph{Ruff documentation}, v0.17: rules, configuration, editor setup.
  \url{https://docs.astral.sh/ruff/}\par
  mypy documentation, v2.4. \url{https://mypy.readthedocs.io}; Microsoft (2026). \emph{Pylance} README and FAQ,
  \url{https://github.com/microsoft/pylance-release}; \emph{Mypy Type Checker}, \url{https://github.com/microsoft/vscode-mypy}\par
\end{frame}

\begin{frame}{A11\enspace References: methods and evidence}
  \scriptsize\setlength\parskip{0.9mm}\vskip-\parskip
  K.~Beck (2002). \emph{Test-Driven Development: By Example}. Addison-Wesley.\par
  M.~Nygard (2011). \emph{Documenting architecture decisions}. Cognitect blog, 15~Nov~2011.
  \url{https://www.cognitect.com/blog/2011/11/15/documenting-architecture-decisions}\par
  N.~S.~Mathews, M.~Nagappan (2024). \emph{Test-driven development for code generation}. ASE~’24.
  \url{https://doi.org/10.1145/3691620.3695527}; arXiv:2402.13521\par
  S.~Kapoor, A.~Narayanan (2023). \emph{Leakage and the reproducibility crisis in machine-learning-based science}.
  Patterns 4(9), 100804. \url{https://doi.org/10.1016/j.patter.2023.100804}\par
  B.-I.~Bachmann, M.~Stiefel, M.~Müller, F.~T.~Mücklich (2026). \emph{A systematically designed open-access SEM
  dataset and robustness benchmark for machine learning-based microstructure segmentation: a high-chromium cast iron
  case study}. Zenodo, v1, CC~BY~4.0. \url{https://doi.org/10.5281/zenodo.21931379}; DFG project MU~959/57-1.\par
\end{frame}

\begin{frame}{A12\enspace References: books and skill packages}
  \scriptsize\setlength\parskip{0.9mm}\vskip-\parskip
  A.~Hunt, D.~Thomas (1999). \emph{The Pragmatic Programmer: From Journeyman to Master}. Addison-Wesley;
  20th anniversary ed., D.~Thomas, A.~Hunt (2019), \emph{The Pragmatic Programmer: Your Journey to Mastery}.\par
  K.~Beck (1999). \emph{Extreme Programming Explained: Embrace Change}. Addison-Wesley; 2nd ed.\ with C.~Andres,
  2004.\par
  M.~Fowler (1999). \emph{Refactoring: Improving the Design of Existing Code}. Addison-Wesley; 2nd ed.\ 2018.
  Chapter~3: bad smells in code.\par
  E.~Evans (2003). \emph{Domain-Driven Design: Tackling Complexity in the Heart of Software}. Addison-Wesley.\par
  J.~Ousterhout (2018). \emph{A Philosophy of Software Design}. Yaknyam Press; 2nd ed.\ 2021.\par
  Agent Skills (2026). \emph{Specification}. \url{https://agentskills.io/specification}; Anthropic (2026). Claude Code
  documentation: \emph{Skills}. \url{https://code.claude.com/docs/en/skills}, both accessed 3~Nov~2026.\par
  J.~Vincent, Prime Radiant (2026). \emph{obra/superpowers}, v7.0.1 (21~Oct~2026), MIT licence.
  \url{https://github.com/obra/superpowers}\par
  J.~Zhang (2026). \emph{andrej-karpathy-skills}, a \fn{CLAUDE.md} derived from A.~Karpathy's observations on LLM coding
  pitfalls. \url{https://github.com/multica-ai/andrej-karpathy-skills}, accessed 3~Nov~2026.\par
  Open GSD (2026). \emph{gsd-core}, v1.20.0 (3~Nov~2026), MIT licence; the continuation of Get Shit Done.
  \url{https://github.com/open-gsd/gsd-core}\par
\end{frame}
